\section*{Факторизация.}

\Def Факторизация - процесс разложения числа делители.

\begin{enumerate}
    \item Линейный алгоритм. Для каждого i от 1 до n проверяем, что число n делиться на i.
\begin{lstlisting}
for (int i = 1; i <= n; ++i) {
  if (n % i == 0) {  
    cout << i << " ";  
  }  
}
\end{lstlisting}
    \item Если $d$ - делитель $n$, то и $n/d$ - делитель, а значит достаточно перебирать не до $n$, а до $[\sqrt{n}]$.
\begin{lstlisting}
for (int i = 1; i <= sqrt(n); ++i) {  
  if (n % i == 0) {  
    cout << i << " ";  
    if (i != n / i) {  
      cout << n / i << " ";  
    }  
  }  
}
\end{lstlisting}
\end{enumerate}

\textbf{Алгоритм разложения числа на простые делители (Оптимизированный).}

\begin{lstlisting}
int d = 2;
while (n > 1) {
  if (n % d == 0) {
    cout << d << " ";
    n /= d;
  } else if (d * d > n) {
    cout << n;
    break;
  } else {
    ++d;
  }
}
\end{lstlisting}

Написан на основе: \href{https://neerc.ifmo.ru/wiki/index.php?title=Разложение_на_множители_(факторизация)#.D0.9F.D1.81.D0.B5.D0.B2.D0.B4.D0.BE.D0.BA.D0.BE.D0.B4_.D0.BD.D0.B0.D1.85.D0.BE.D0.B6.D0.B4.D0.B5.D0.BD.D0.B8.D1.8F_.D0.BF.D1.80.D0.BE.D1.81.D1.82.D1.8B.D1.85_.D0.BC.D0.BD.D0.BE.D0.B6.D0.B8.D1.82.D0.B5.D0.BB.D0.B5.D0.B9}{Викиконспекты (ИТМО). Псевдокод нахождения простых делителей.}

\section*{Функция Эйлера.}

\Def Функция Эйлера от числа $n, \varphi(n)$, равна количеству чисел, меньших $n$ и взаимно простых с ним. 

Например, $\varphi(5) = 4$ (числа $1, 2, 3, 4$).

\textbf{Свойства функции Эйлера.}

\begin{enumerate}
    \item Если $p$ - простое, то $\varphi(p) = p - 1$.
    \item Если $p$ - простое, то $\varphi(p^\alpha) = p^\alpha - p^{\alpha - 1}$.
    \item Если $n, m$ взаимно просты, то $\varphi(m \cdot n) = \varphi(m) \cdot \varphi(n)$. 
\end{enumerate}

\textbf{Сводимость вычисления функции Эйлера к факторизации.}

Наивный алгоритм. Для каждого числа от 1 до $n$ проверяем, является ли делителем $n$.

Немного преобразований:

$$\varphi \brackets{n} = \varphi \brackets{\prod_{i}{p_{i}^{\alpha_i}}} = \prod_{i}\varphi\brackets{p^{\alpha_i}_{i}} = \prod_{i}\brackets{p^{\alpha_i}_{i} - p^{\alpha_i - 1}_{i}} = $$

$$\prod_{i}\brackets{p^{\alpha_i}_{i}\brackets{1 - \frac{1}{p_i}}}=\prod_{i}p^{\alpha_i}_{i} \cdot \prod_i\brackets{1 - \frac{1}{p_i}} = n \cdot \prod_i\brackets{1 - \frac{1}{p_i}}$$

Таким образом, зная факторизацию числа, можно вычислить $\varphi(n)$ за $\mathcal{O}(\log n)$.

\pagebreak